\documentclass{article}
\usepackage{graphicx,color}
\usepackage{parskip}
\begin{document}

Jan, Joachim has run a bunch of simulations and we have looked through them in order to find a reasonable story. 
Joachim, I already lost the sketch due to moving between offices etc. In planned to write this in google
docs, but was not able to upload the pdfs as images.  

Anyway, here goes the story. 


\begin{figure}
\begin{center}
\includegraphics[page=24,width=\linewidth]{presentation.pdf}
\caption{2D spinal cord problem. The parameters are as shown in the figure, except that $b=0$ everywhere (giving $\tau\rightarrow\infty$), and $\Delta t=10^{-3}$.}
\end{center}
\end{figure}

\begin{figure}
\begin{center}
\includegraphics[width=7cm]{Richardson,problem=2,exact=1,N=64.pdf}
\caption{Richardson iterations (exact subsolves) is equivalent to the traditional splitting schemes. The residual is shown as solid lines; energy-norm error as dotted lines. Coupling strength $\tau\rightarrow\infty$ when $b\rightarrow 0$, so the unstable splittings diverge quickly. (Undrained split is not included in the comparison, because it is undefined for $b=0$.)} 
\label{2dsc-richardson-exact}
\end{center}
\end{figure}

\begin{figure}
\begin{center}
\includegraphics[width=7cm]{BiCGStab,problem=2,exact=1,N=64.pdf}
\caption{BiCGStab (exact subsolves). The residuals is shown as solid lines; energy-norm error as dotted lines.}
\label{2dsc-bicgstab-exact}
\end{center}
\end{figure}

\begin{figure}
\begin{center}
\includegraphics[width=7cm]{LGMRES,problem=2,exact=1,N=64,cycles=1.pdf}
\caption{LGMRES (exact subsolves). The residuals is shown as solid lines; energy-norm error as dotted lines. The latter are (nearly?) invisible because they overlap with the residual lines.}
\label{2dsc-lgmres-exact}
\end{center}
\end{figure}

Pictures \ref{2dsc-richardson-exact}--\ref{2dsc-lgmres-exact} shows that Krylov methods are superior to Richardson. This is not a big surprise, but provides
a clear link to the splitting schemes.

\begin{figure}
\begin{center}
\includegraphics[width=7cm]{Richardson,problem=2,exact=0,N=64.pdf}
\includegraphics[width=7cm]{LGMRES,problem=2,exact=0,N=64,cycles=1.pdf}
\caption{Richardson and LGMRES with inexact subsolves (1 cycle of ML).} 
\label{2dsc-richardson-lgmres-inexact}
\end{center}
\end{figure}

\begin{figure}
\begin{center}
\includegraphics[width=7cm]{BiCGStab,problem=2,exact=0,N=64.pdf}
\caption{BiCGStab} 
\end{center}
\end{figure}

\begin{figure}
\begin{center}
\includegraphics[width=7cm]{LGMRES,problem=2,exact=0,N=64,cycles=1.pdf}
\caption{LGMRES} 
\end{center}
\end{figure}

The figures also show both the error (in energy norm) and the residual for BiCGStab and LGMRES. 
Here stripled lines are residual and the solid line is energy. Here, it is clear that it is a
close correspondance between error and residual for LGMRES and less clear for BiCGStab. We therefore
continue the discussion with only LGMRES and error.


\emph{JBH: It is not immediately obvious that the Krylov methods are very
  superior, given that BiCGStab is about twice as expensive per iteration, and
  LGMRES is 30+ times as expensive per iteration. What is clear is: (a) They
  are reasonable preconditioners, compared to the bog-standard Block Jacobi
  method; (b) Even the unstable splittings converge, at least in LGMRES; and
  (c) Error is much better controlled in LGMRES than in either Richardson or
  BiCGStab.}

Figure \ref{2dsc-richardson-lgmres-inexact} shows that the same pattern holds
when the splitting is performed using inexact single-block solvers.

Next, lets consider the eigenvalues of the individual components. 
\begin{figure}
\begin{center}
\includegraphics[width=7cm]{EV[A],problem=2,exact=0,N=64.pdf}
\caption{Eigenvalues} 
\end{center}
\end{figure}

\begin{figure}
\begin{center}
\includegraphics[width=7cm]{EV[-C],problem=2,exact=0,N=64.pdf}
\caption{Eigenvalues} 
\end{center}
\end{figure}


\begin{figure}
\begin{center}
\includegraphics[width=7cm]{EV[-S_c],problem=2,exact=0,N=64.pdf}
\caption{Eigenvalues} 
\end{center}
\end{figure}

Clearly, the eigenvalues here are nicely bounded. We may here investigate how 
the number of ML sweeps affect the eigenvalues.  

This more or less concludes the discussion of the 2D spinal cord example. 

\emph{JBH: Can we drop everything except LGMRES on (inexact) fixed stress in the rest of the examples?}

\begin{figure}
\begin{center}
\includegraphics[page=36,width=\linewidth]{presentation.pdf}
\caption{Contrast problem. There are differences in the parameters. Again $b=0$ everywhere (giving $\tau\rightarrow\infty$), and $\Delta t=10^{-3}$. Also, $p=0$ at the top, and $E$ is 1/10 of that shown here. These changes are arbitrary: I can create new figures with the same parameters if necessary.}
\end{center}
\end{figure}

\begin{figure}
\begin{center}
\includegraphics[width=7cm]{Richardson,problem=1,exact=0,N=20.pdf}
\caption{Contrast problem, Richardson, $N=20$.}
\end{center}
\end{figure}

\begin{figure}
\begin{center}
\includegraphics[width=7cm]{Richardson,problem=1,exact=0,N=160.pdf}
\caption{Contrast problem, Richardson, $N=160$.}
\end{center}
\end{figure}

\begin{figure}
\begin{center}
\includegraphics[width=7cm]{LGMRES,problem=1,exact=0,N=20,cycles=1.pdf}
\caption{Contrast problem, LGMRES, $N=20$.}
\end{center}
\end{figure}

\begin{figure}
\begin{center}
\includegraphics[width=7cm]{LGMRES,problem=1,exact=0,N=160,cycles=1.pdf}
\caption{Contrast problem, LGMRES, $N=160$.}
\end{center}
\end{figure}


\end{document}

